\documentclass[landscape,8pt]{article}
\linespread{1}
\usepackage{multicol}
\setlength{\columnsep}{0.5cm}
\setlength{\columnseprule}{0.25pt}
\usepackage{amsmath}
\usepackage{graphbox}
\usepackage{amssymb}
\usepackage{amsfonts}
\usepackage{mathtools}
\usepackage{subcaption}
\usepackage{geometry}
\geometry{a4paper, landscape, margin = 0.6in}
\usepackage{enumerate}
\usepackage{amsthm}
\usepackage{physics}
\usepackage{float}
\usepackage{kotex}
\usepackage{tabu}
\usepackage{listings}
\usepackage{wrapfig}
\usepackage{authblk}
\usepackage{verbatim}
\usepackage{graphicx}
\usepackage{xcolor}
\usepackage{fancyhdr}

\pagestyle{fancy}
\fancyhead[R]{Little Piplup}
\fancyhead[L]{Seoul National University}
\usepackage[compact]{titlesec}
\DeclarePairedDelimiter{\ceil}{\lceil}{\rceil}
\DeclarePairedDelimiter{\floor}{\lfloor}{\rfloor}
\newcommand{\st}{\text{ such that }}
\newcommand{\for}{\text{ for }}
\newcommand{\newpara}[1]{\paragraph{#1} \mbox{}\\}
\everymath{\displaystyle}
\setlength{\parskip}{0pt}
\setlength{\parsep}{0pt}
\setlength{\parindent}{0pt}
\titlespacing{\section}{0pt}{*0}{*0}
\titlespacing{\subsection}{0pt}{*1}{*0.5}
\titlespacing{\subsubsection}{0pt}{*0}{*0}
\topsep=5pt \partopsep=5pt
\title{\Huge{\textsf{Little Piplup}}}
\author{Gratus907(Wonseok Shin), Coffeetea(Jinje Han), DhDroid(Donghyun Son)}
\date{ }

\begin{document}
{
  \setlength{\columnseprule}{0pt}
  \begin{multicols}{3}\raggedcolumns
  \vspace*{\fill}
  \begin{figure}[H]
  \includegraphics[width=\linewidth]{piplup}
  \end{figure}
  \vspace*{\fill}
  \columnbreak
  \begin{center}
  \vspace*{\fill}
  \huge{\textbf{Little Piplup}\\}
  \vspace{2em}
    \LARGE{Contest Teamnote \\}
   \large{UCPC 2019 ver}
  \vspace*{\fill}
  \end{center}
  \columnbreak

  \vspace*{\fill}
  \begin{figure}[H]
  \includegraphics[width=\linewidth]{revplup}
  \end{figure}

  \vspace*{\fill}
  \end{multicols}
}




\newpage
\maketitle
\begin{multicols}{3}\raggedcolumns
\tableofcontents
\end{multicols}
\pagebreak
\begin{multicols}{2}\raggedcolumns


\section{Settings}

  \subsection{C++}

\section{Data Structures}
  \subsection{Segment Tree}
     To deal with queries on intervals, we use segment tree.
     \verbatiminput{./code/Data_Structures/Segtree_Make.cpp}
     Updating segment tree
     \verbatiminput{./code/Data_Structures/Segtree_Update.cpp}
     Answering queries with segment tree
     \verbatiminput{./code/Data_Structures/Segtree_Answerquery.cpp}
     Answering range minimum queries with segment tree
     \verbatiminput{./code/Data_Structures/Range_Minimum_Query.cpp}
    \columnbreak
  %\subsection{Fenwick Tree}
  \subsection{Disjoint Set Union (Union - Find)}
      \verbatiminput{./code/Data_Structures/Disjoint_Set_Union.cpp}

\columnbreak
\section{Mathematics}
  \subsection{Useful Mathematical Formula}
    \begin{itemize}
      \item Catalan Number : Number of valid parantheses strings with $n$ pairs
      \[
        C_n = \frac{1}{n+1}\binom{2n}{n}
      \]
      \item Nim Game : Remember - XOR of all piles.
    \end{itemize}
  \subsection{Number of Integer Partition}
    \verbatiminput{./code/Mathematics/Num_of_Integer_Partition.py}

  \subsection{Binomial Coefficient}
    Fast-to-Type Binomial coefficient

  \subsection{Extended Euclidean Algorithm}
    \verbatiminput{./code/Mathematics/Extended_Euclidean_Algorithm.cpp}

  \subsection{Fast Modulo Exponentiation}
  Calculating \texttt{$x^y$ mod $p$} in $\order{\log y}$ time.
    \verbatiminput{./code/Mathematics/Fast_Modulo_Exponentiation.cpp}

  \subsection{Miller-Rabin Primality Testing}
  Base values of $a$ chosen so that results are tested to be correct up to $10^14$.
    \verbatiminput{./code/Mathematics/Miller_Rabin_Test.cpp}

  \subsection{Pollard-Rho Factorization}
    \verbatiminput{./code/Mathematics/Pollard_Rho_Algorithm.cpp}

  \subsection{Euler Totient}
  Calculating number of integers below $n$ which is coprime with $n$.
    \verbatiminput{./code/Mathematics/Euler_Phi_Function.cpp}

  \subsection{Modular Multiplicative Inverse}
    \verbatiminput{./code/Mathematics/Modular_Inverse.cpp}

  %\subsection{Kitamasa method}

  %\subsection{Fast Fourier Transform}
  \columnbreak
\columnbreak
\section{Geometry}
  \subsection{CCW}
    \verbatiminput{./code/Geometry/CCW.cpp}

  \subsection{Point in polygon}
    Returns boolean, if point is in the polygon (represented as vector of points).
    \verbatiminput{./code/Geometry/Point_in_polygon.cpp}

  \subsection{Closest Pair Problem}
    Requires : Points must be sorted with x-axis.\\
    Runs in $\mathcal{O}(n \log^2 n)$
    \verbatiminput{./code/Geometry/Closest_pair.cpp}
  \subsection{Convex Hull (Graham Scan)}
    \verbatiminput{./code/Geometry/Graham_Scan.cpp}
  \subsection{Intersection of Line Segment}
    \verbatiminput{./code/Geometry/Line_Segment_Intersection.cpp}
\columnbreak
\section{Graphs}
  \subsection{Topological Sorting}
  Topological sorting with dfs
    \verbatiminput{./code/Graph/Topological_Sort.cpp}

  \subsection{Lowest Common Ancestor}
  LCA Algorithm by sparse table
    \verbatiminput{./code/Graph/Lowest_Common_Ancestor.cpp}

  \subsection{MST Kruskal Algorithm}
  Based on Union-Find implementation\\
  $\order{E \log E}$ if path-compressed Union Find.
    \verbatiminput{./code/Graph/Kruskal_MST.cpp}

  \subsection{MST Prim Algorithm}
    \verbatiminput{./code/Graph/Prim_MST.cpp}

  %\subsection{MST Borvuka Algorithm}

  \subsection{Dinic's Algorithm}
    \verbatiminput{./code/Graph/Dinic.cpp}
  %\subsection{Heavy-Light Decomposition}

  %\subsection{Centroid Decomposition}

  %\subsection{Hungarian Algorithm}
    %$\order{n^3}$ assignment problem
\columnbreak

\section{Shortest Path}
  \subsection{Dijkstra}
    $\order{E \log V}$ Single-Start-Shortest-Path.\\
    Not working for graph with minus weight.
      \verbatiminput{./code/Graph/Dijkstra_Algorithm.cpp}

  \subsection{Bellman Ford}
    $\order{EV}$ Single-Start-Shortest-Path.\\
    Not working for graph with minus cycle $\rightarrow$ must detect.
      \verbatiminput{./code/Graph/Bellman_Ford_Algorithm.cpp}

  \subsection{SPFA Algorithm}
    Average $\order{E}$, worst $\order{VE}$ time. Average-case improvement of Bellman Ford by using an additional queue. $\rightarrow$ 데이터를 누가 짰을지를 생각해 보면 그냥 이런거 집어치우는게 맞을듯. ICPC Preliminary에서나 쓰자.

  \subsection{Floyd-Warshall}
    Works on adjacency matrix, in $\order{V^3}$.
      \verbatiminput{./code/Graph/Floyd_Warshall_Algorithm.cpp}

\columnbreak

\section{Dynamic}
  \subsection{Longest Increasing Subsequence}
    Find LIS in $\order{n \log n}$ time.
      \verbatiminput{./code/Dynamic_Programming/Longest_Increasing_Subsequence.cpp}

  \subsection{Largest Sum Subarray}
    Computes sum of largest sum subarray in $\order{N}$
      \verbatiminput{./code/Dynamic_Programming/Largest_Maximum_Subarray.cpp}

  \subsection{0-1 Knapsack}
    \verbatiminput{./code/Dynamic_Programming/0-1_Knapsack.cpp}
  \subsection{Longest Common Subsequence}
    \verbatiminput{./code/Dynamic_Programming/Longest_Common_Substring.cpp}

  %\subsection{Edit Distance}

  %\subsection{Convex Hull Trick}

  %\subsection{Divide and Conquer Optimization}

  %\subsection{Knuth Optimization}

\section{String}
  \subsection{KMP Algorithm}
    \verbatiminput{./code/String/KMP.cpp}
  \subsection{Manacher's Algorithm}
    \verbatiminput{./code/String/Manacher.cpp}
  \subsection{Trie}
    \verbatiminput{./code/String/Trie.cpp}

  %\subsection{Rabin-Karp Hashing}

  %\subsection{Aho-Corasick Algorithm}
\columnbreak

\section{Miscellaneous}
  \subsection{Useful Bitwise Functions in C++}
  \begin{verbatim}
    int __builtin_clz(int x);// number of leading zero
    int __builtin_ctz(int x);// number of trailing zero
    int __builtin_clzll(ll x);// number of leading zero
    int __builtin_ctzll(ll x);// number of trailing zero
    int __builtin_popcount(int x);// number of 1-bits in x
    int __builtin_popcountll(ll x);// number of 1-bits in x

    lsb(n): (n & -n); // last bit (smallest)
    floor(log2(n)): 31 - __builtin_clz(n | 1);
    floor(log2(n)): 63 - __builtin_clzll(n | 1);

    // compute next perm. ex) 00111, 01011, 01101, 01110, 10011, 10101..
    ll next_perm(ll v)
    {
      ll t = v | (v-1);
      return (t + 1) | (((~t & -~t) - 1) >> (__builtin_ctz(v) + 1));
    }
  \end{verbatim}
  \subsection{List of Useful Numbers}
  \begin{verbatim}

    < 10^k    prime   # of prime          < 10^k            prime
    -------------------------------------------------------------
    1             7            4          10           9999999967
    2            97           25          11          99999999977
    3           997          168          12         999999999989
    4          9973         1229          13        9999999999971
    5         99991         9592          14       99999999999973
    6        999983        78498          15      999999999999989
    7       9999991       664579          16     9999999999999937
    8      99999989      5761455          17    99999999999999997
    9     999999937     50847534          18   999999999999999989
  \end{verbatim}
\columnbreak


\section{Checkpoints}
  \subsection{Debugging}
\begin{itemize}
  \item $10^5 * 10^5 \Rightarrow \text{INTEGER OVERFLOW}$.
  \item If unsure with overflow, use \\
  \texttt{\#define int long long} and stop caring.
  \item 행렬과 기하의 $i, j$ 인덱스 조심. 헷갈리면 쓰면서 가기.
  \item output이 특정 수열/OX 형태 : 작은 예제를 Exhasutive Search. 모르는 무언가를 알기 위해서는 데이터가 필요하다.
\end{itemize}
  \subsection{Thinking}
    \begin{itemize}
      \item 모든 경우를 다 할 수 없나? 왜 안 되지? 시간 복잡도 잘 생각해 보기. 정해의 Target Complexity를 먼저 생각하고 주요 알고리즘들의 Complexity로 짜맞추기.\\
      예를들어, 쿼리가 30만개 들어온다면 한 쿼리를 적어도 $\log{n}$ 에 처리할 방법이 아무튼 있다는 뜻.
      \item 단조함수이며, 충분히 빠르게 검증가능한가 : Binary Search.
      \item 차원이 높은 문제 : 차원 내려서 생각하기. 3 $\rightarrow$ 2.
      \item 이 문제가 사실 그래프 관련 문제는 아닐까?
        \begin{itemize}
          \item 만약 그렇다면, `간선' 과 `정점' 은 각각..?
        \end{itemize}
      \item 이 문제에 Overlapping Subproblem이 보이나? $\rightarrow$ Dynamic Programming 을 적용.
      \item 답의 상한이 Reasonable 하게 작은가?
      \item 답이 단순한 수열이면, Berlekamp-Massey (본선에서) or OEIS (예선)
      \item 마지막 생각 : 조금 추하지만 해싱이나 Random \texttt{bitset} 을 이용한 $n^2 / 64$ 같은걸로 뚫을 수 있나?
    \end{itemize}
\end{multicols}
\end{document}
